\subsection{Useful wall information remains in the frozen representation}
\label{sec:results-wall}

Decoded real encodings visibly retain the approaching wall (Figure~\ref{fig:rollouts}). To test
whether that information can support behavior, we train a convolutional policy in the real gate
on the frozen decoder's output. Two of nine runs discover the survival strategy, both using encoder
seed 1; the corresponding real-pixel CNN discovers it in two of three runs
(Figure~\ref{fig:representation}). This establishes useful information in at least one fitted encoder
under the decoder-and-CNN read-out. The decoder adds reconstruction supervision outside the
world-model budget, so this is a test of recoverable information rather than unsupervised control
sufficiency.

The policy interface matters. With an MLP torso, real pixels yield $0.91$ walls per episode and
frozen latents yield $0.94$, with no discoveries in either group. An MLP on decoded latents similarly
yields $0.92$ walls. Thus failure of this MLP training recipe is not evidence that the encoder has
lost the obstacle. These comparisons do not set a capacity limit for all MLP policies; they locate
a limitation of the tested read-out and optimization procedure.

\begin{figure}[t]
\centering
\paperfig{rollouts.pdf}{1.0}
\caption{\textbf{A wall visible in real encodings disappears in predicted rollouts.} A pre-reveal
context is followed by futures under recorded actions: the simulator, decoded real JEPA encodings,
decoded original JEPA predictions, and one flow-model sample. The JEPA rows use encoder seed 0
and the same fitted decoder. The example is the first stored diagnostic context at wall position
$x=45$ among 64 contexts whose recorded futures survive the horizon. The flow sample is a single-checkpoint feasibility reference.
A decoding failure on predicted latents can reflect both prediction error and off-distribution
read-out; it is not alone proof of information loss.}
\label{fig:rollouts}
\end{figure}

\subsection{The tested imagination pipeline does not discover the strategy}
\label{sec:results-ppo}
\label{sec:results}

All three JEPA predictor variants produce poor downstream control
(Table~\ref{tab:gate-main}; per-seed outcomes in Appendix~\ref{app:round2}). The original predictor yields $0.27$ walls
per episode; one residual sample and the mean of eight each yield $0.26$. No policy among these
27 runs discovers the strategy; the highest final survival is $3.1\%$. Real-transition PPO on the
same frozen encoders yields $0.94$ walls. The difference measures the combined effect of learned
transitions, fitted heads, and the changed PPO protocol. Even the stronger real-latent MLP result
falls far short of reliable survival.

The flow reference supports feasibility at this task and nominal model budget: one of its nine
policies discovers the strategy. The remaining eight pass $0.72$--$1.18$ walls per episode. Since
all nine use one dynamics fit with a different architecture and supervision, this reference does
not establish a model-family ranking. Training the gate JEPA for twice and four times the steps
also leaves imagined-policy survival at $0$--$3\%$ in the tested extensions (encoder seed 0, policy
seed 0, each predictor variant).

\begin{table}[t]
\centering
\caption{\textbf{Final gate outcomes, with predictor variants kept separate.} Mean $\pm$ sample
standard deviation across policy runs; discovery is final survival above $50\%$. JEPA imagination variants and real-latent
conditions each contain $3$ models $\times$ $3$ policies; pixel controls have three policies. Real training changes rewards, terminations, and PPO protocol
alongside dynamics. Flow is a single-model reference.}
\label{tab:gate-main}
\inputtable{gate_main.tex}
\end{table}



\subsection{Prediction and fitted heads provide distinct failure signals}
\label{sec:results-mean}

The same decoder that renders real encoded walls removes them from original predicted rollouts.
Across all three encoder seeds, decoded shared-wall intensity is at most $0.02$ by the second
predicted step (Appendix~\ref{app:round2}). The shared wall pixels occur in both possible gap futures,
so a pixelwise mixture of those futures would retain them. This observation rules out that simple
pixel-space picture; it does not determine what averaging or distribution shift does in the learned
latent space.

A within-encoder counterfactual diagnostic makes the scale of the feature explicit. On 64 pre-reveal
contexts for seed 0, the first-step squared distance between wall-present and wall-absent encodings
is $0.035$, while total teacher-forced prediction error is $0.145$. The same ordering holds for all
three seeds (Appendix~\ref{app:round2}). This comparison concerns sensitivity to wall removal versus
error in all features. It neither identifies a causal mechanism nor permits comparing information
content through raw distances across independently fitted latent spaces.

The fitted heads expose another limitation. For encoder 0, termination AUC is $0.96$ for border
deaths but $0.65$ for wall deaths. Predicted rewards for the two death causes are $-0.89$ and $-0.05$,
respectively, against the true $-1$. An imagined policy therefore trains on a weak wall-collision
signal. The residual-sampling variants retain these heads and the policy protocol. Their poor control
outcomes therefore test this particular intervention, not distributional prediction in general.

\begin{figure}[t]
\centering
\paperfig{representation_by_seed.pdf}{1.0}
\caption{\textbf{Recoverable information can support successful control.} Final real-gate survival
for MLP policies on frozen latents or pixels (left), and CNN policies on decoded latents or pixels
(right). Points are policy runs, grouped by encoder seed; the dashed line is the $50\%$ discovery criterion. Both decoded-latent discoveries
use encoder 1. Reconstruction supervision and CNN training are additional to representation fitting.}
\label{fig:representation}
\end{figure}
